% META: {"id": "1NSI-AGT-COURS-C2", "chapitre": "1NSI-ALGO-PARCOURS-TRIS", "type_objet": "cours", "capacites": ["P-ALGO-02A", "P-ALGO-02C"], "mode_creation": "ex_nihilo", "sources_inspiration": [], "status": "generated"}
\section{Le tri par insertion}

\definition[D2]{
  Le \textbf{tri par insertion} construit progressivement une partie triée du tableau, en
  \textbf{insérant}, un par un, chaque élément restant à sa bonne place dans la partie déjà
  triée (comme on trie des cartes à jouer en main, une par une).
}

\begin{codereference}{Tri par insertion}
\begin{python}
def tri_insertion(tableau):
    """Trie `tableau` par ordre croissant, par insertion (renvoie une copie triee)."""
    tableau = tableau[:]
    for i in range(1, len(tableau)):
        valeur = tableau[i]
        j = i - 1
        while j >= 0 and tableau[j] > valeur:
            tableau[j + 1] = tableau[j]
            j -= 1
        tableau[j + 1] = valeur
    return tableau

print(tri_insertion([5, 3, 8, 1, 9, 2]))  # [1, 2, 3, 5, 8, 9]
\end{python}
\end{codereference}

% BEGIN-VERIFY
% def tri_insertion(tableau):
%     tableau = tableau[:]
%     for i in range(1, len(tableau)):
%         valeur = tableau[i]
%         j = i - 1
%         while j >= 0 and tableau[j] > valeur:
%             tableau[j + 1] = tableau[j]
%             j -= 1
%         tableau[j + 1] = valeur
%     return tableau
% assert tri_insertion([5, 3, 8, 1, 9, 2]) == [1, 2, 3, 5, 8, 9]
% assert tri_insertion([]) == []
% assert tri_insertion([1]) == [1]
% assert tri_insertion([2, 1]) == [1, 2]
% END-VERIFY

\propriete[Terminaison et coût]{
  La boucle \lstinline{for} externe s'exécute un nombre \textbf{fixe} de fois
  ($n-1$ itérations pour un tableau de taille $n$) : elle termine nécessairement. La boucle
  \lstinline{while} interne termine car $j$ \textbf{décroît strictement} à chaque tour et
  est bornée par $0$. Dans le \textbf{pire cas} (tableau trié en ordre décroissant), le
  coût est \textbf{quadratique} ($O(n^2)$) : chaque insertion peut nécessiter de décaler
  tous les éléments déjà triés.
}

\demonstration{
\textbf{Invariant de boucle} : au début de chaque itération $i$ de la boucle \lstinline{for}
(avant l'insertion de \lstinline{tableau[i]}), le sous-tableau \lstinline{tableau[0:i]} est
\textbf{trié} (et contient les mêmes éléments que le sous-tableau original
\lstinline{tableau[0:i]} avant tout tri).

\textbf{Initialisation.} Avant la première itération ($i=1$), \lstinline{tableau[0:1]} est
un tableau à un seul élément : il est trivialement trié.

\textbf{Conservation.} Supposons \lstinline{tableau[0:i]} trié au début de l'itération $i$.
La boucle \lstinline{while} interne décale vers la droite tous les éléments de
\lstinline{tableau[0:i]} strictement supérieurs à \lstinline{valeur = tableau[i]}, puis
place \lstinline{valeur} à la position ainsi libérée. Le sous-tableau
\lstinline{tableau[0:i+1]} obtenu est donc trié : l'invariant est conservé pour
l'itération suivante.

\textbf{Terminaison.} Après la dernière itération ($i=n-1$), l'invariant donne :
\lstinline{tableau[0:n]} (le tableau entier) est trié. L'algorithme est donc
\textbf{correct}.
}

% BEGIN-VERIFY
% def tri_insertion_avec_invariant(tableau):
%     tableau = tableau[:]
%     for i in range(1, len(tableau)):
%         # invariant avant insertion : tableau[0:i] est trie
%         assert tableau[0:i] == sorted(tableau[0:i])
%         valeur = tableau[i]
%         j = i - 1
%         while j >= 0 and tableau[j] > valeur:
%             tableau[j + 1] = tableau[j]
%             j -= 1
%         tableau[j + 1] = valeur
%         # invariant conserve apres insertion : tableau[0:i+1] est trie
%         assert tableau[0:i + 1] == sorted(tableau[0:i + 1])
%     return tableau
% resultat = tri_insertion_avec_invariant([5, 3, 8, 1, 9, 2])
% assert resultat == [1, 2, 3, 5, 8, 9]
% END-VERIFY

\erreurFrequente{Croire qu'un invariant de boucle doit être vérifié \textbf{une seule fois}
(par exemple seulement à la fin). Un invariant doit être vrai \textbf{à chaque} itération :
c'est cette répétition, de l'initialisation jusqu'à la terminaison, qui constitue la preuve
par récurrence de la correction de l'algorithme.}

\alamachine{Ajoute des instructions \lstinline{assert tableau[0:i] == sorted(tableau[0:i])}
dans ta propre implémentation du tri par insertion, juste avant l'insertion de chaque
nouvel élément, pour vérifier expérimentalement l'invariant sur plusieurs tableaux de ton
choix.}
